\textbf{\textit{Activity 7. }Computer y submit the exercises from 24 to 27 of the page 60 of [2]}
\begin{quote}

\textbf{Bernoulli Equations.} Sometimes it is possible to solve a nonlinear equation by making a change of the dependent variable that converts it into a linear equation. The most important such equation has the form
\[
y' + p(t)y = q(t)y^n,
\]
and is called a Bernoulli equation after Jakob Bernoulli. Problems 23 and 25 deal with equations of this type.


\textbf{24. $y' = ry - ky^2$, $r > 0$ and $k > 0$. This equation is important in population dynamics and is discussed in detail in Section 2.5.}
\begin{respuesta}

\begin{theorem}
This part was so fucking hard, so I used heavy IA to solve it, as well the time for this assignment wasn't enough, I didn't put enough time for learning the real purpose of the "Bernoulli" stuff, good reading anyway. 
\end{theorem}
\vspace{0.5cm}


This is a Bernoulli equation with $p(t) = -r$, $q(t) = -k$, and $n = 2$.

\textbf{Rewrite in standard Bernoulli form.}

\[
y' - ry = -ky^2
\]

\textbf{Apply the Bernoulli substitution.}

Divide both sides by $y^n = y^2$:

\[
y^{-2}y' - ry^{-1} = -k
\]

Let $v = y^{1-n} = y^{-1}$, so that $v' = -y^{-2}y'$, which gives $y^{-2}y' = -v'$.

Substituting:

\[
-v' - rv = -k
\]

\[
v' + rv = k
\]

\textbf{Solve the linear equation.}

This is a first-order linear ODE in $v$. The integrating factor is:

\[
\mu(t) = e^{\int r\, dt} = e^{rt}
\]

Multiplying both sides by $\mu(t)$:

\[
\frac{d}{dt}\left(e^{rt}v\right) = ke^{rt}
\]

Integrating both sides:

\[
e^{rt}v = \frac{k}{r}e^{rt} + C
\]

\[
v = \frac{k}{r} + Ce^{-rt}
\]

\textbf{Back-substitute $v = y^{-1}$.}

\[
\frac{1}{y} = \frac{k}{r} + Ce^{-rt}
\]

\[
\boxed{y = \frac{1}{\dfrac{k}{r} + Ce^{-rt}}}
\]

\textbf{Analyze the behavior.}

As $t \to \infty$, since $r > 0$ we have $e^{-rt} \to 0$, so:

\[
y \to \frac{1}{k/r} = \frac{r}{k}
\]

This confirms that the population approaches the carrying capacity $\dfrac{r}{k}$ regardless of the initial condition (provided $y(0) > 0$), which is the hallmark of logistic growth.

\end{respuesta}



\textbf{25. $y' = \epsilon y - \sigma y^3$, $\epsilon > 0$ and $\sigma > 0$.}
\begin{respuesta}

This is a Bernoulli equation with $p(t) = -\epsilon$, $q(t) = -\sigma$, and $n = 3$.

\textbf{Standard Bernoulli form.}

\[
y' - \epsilon y = -\sigma y^3
\]

\textbf{Bernoulli substitution.}

Divide both sides by $y^n = y^3$:

\[
y^{-3}y' - \epsilon y^{-2} = -\sigma
\]

Let $v = y^{1-n} = y^{-2}$, so that $v' = -2y^{-3}y'$, which gives $y^{-3}y' = -\dfrac{v'}{2}$.

Substituting:

\[
-\frac{v'}{2} - \epsilon v = -\sigma
\]

Multiplying through by $-2$:

\[
v' + 2\epsilon v = 2\sigma
\]

\textbf{Solving the linear equation.}

This is a first-order linear ODE in $v$. The integrating factor is:

\[
\mu(t) = e^{\int 2\epsilon\, dt} = e^{2\epsilon t}
\]

Multiplying both sides by $\mu(t)$:

\[
\frac{d}{dt}\left(e^{2\epsilon t}v\right) = 2\sigma e^{2\epsilon t}
\]

Integrating both sides:

\[
e^{2\epsilon t}v = \frac{\sigma}{\epsilon}e^{2\epsilon t} + C
\]

\[
v = \frac{\sigma}{\epsilon} + Ce^{-2\epsilon t}
\]

\textbf{Back-substitution $v = y^{-2}$.}

\[
\frac{1}{y^2} = \frac{\sigma}{\epsilon} + Ce^{-2\epsilon t}
\]

\[
\boxed{y = \frac{1}{\sqrt{\dfrac{\sigma}{\epsilon} + Ce^{-2\epsilon t}}}}
\]

\textbf{Long-term behavior.}

As $t \to \infty$, since $\epsilon > 0$ we have $e^{-2\epsilon t} \to 0$, so:

\[
y \to \frac{1}{\sqrt{\sigma/\epsilon}} = \sqrt{\frac{\epsilon}{\sigma}}
\]

Thus the solution approaches the equilibrium value $\sqrt{\epsilon/\sigma}$ as $t \to \infty$,
indicating that small perturbations in the fluid flow decay toward this stable state.

\end{respuesta}

\noindent\textbf{Discontinuous Coefficients.} Linear differential equations sometimes occur in which one or both of the functions $p$ and $g$ have jump discontinuities. If $t_0$ is such a point of discontinuity, then it is necessary to solve the equation separately for $t < t_0$ and $t > t_0$. Afterward, the two solutions are matched so that $y$ is continuous at $t_0$; this is accomplished by a proper choice of the arbitrary constants. The following two problems illustrate this situation. Note in each case that it is impossible also to make $y'$ continuous at $t_0$.



\textbf{26. Solve the initial value problem
    \[ y' + 2y = g(t), \qquad y(0) = 0, \]
    where
    \[ g(t) = \begin{cases} 1, & 0 \le t \le 1, \\ 0, & t > 1. \end{cases} \]}
\begin{respuesta}

The function $g(t)$ has a jump discontinuity at $t_0 = 1$, so we solve the equation
separately on $[0,1]$ and $(1, \infty)$, then match the solutions by requiring $y$ to be
continuous at $t = 1$.

\textbf{Integrating factor.}

For both intervals the left-hand side has the same form. The integrating factor is:

\[
\mu(t) = e^{\int 2\, dt} = e^{2t}
\]

so the equation becomes in each interval:

\[
\frac{d}{dt}\left(e^{2t}y\right) = g(t)e^{2t}
\]

\textbf{Solution on $0 \le t \le 1$.}

Here $g(t) = 1$, so:

\[
\frac{d}{dt}\left(e^{2t}y\right) = e^{2t}
\]

\[
e^{2t}y = \frac{1}{2}e^{2t} + C_1
\]

\[
y_1 = \frac{1}{2} + C_1 e^{-2t}
\]

Applying $y(0) = 0$:

\[
0 = \frac{1}{2} + C_1 \implies C_1 = -\frac{1}{2}
\]

\[
y_1 = \frac{1}{2}\left(1 - e^{-2t}\right), \qquad 0 \le t \le 1
\]

\textbf{Solution on $t > 1$.}

Here $g(t) = 0$, so:

\[
\frac{d}{dt}\left(e^{2t}y\right) = 0
\]

\[
e^{2t}y = C_2
\]

\[
y_2 = C_2 e^{-2t}, \qquad t > 1
\]

\textbf{Matching at $t = 1$.}

We require $y$ to be continuous at $t = 1$, i.e.\ $y_1(1) = y_2(1)$:

\[
\frac{1}{2}\left(1 - e^{-2}\right) = C_2 e^{-2}
\]

\[
C_2 = \frac{e^{2} - 1}{2}
\]

\textbf{Complete solution.}

\[
\boxed{y = \begin{cases} \dfrac{1}{2}\!\left(1 - e^{-2t}\right), & 0 \le t \le 1, \\[8pt] \dfrac{e^{2}-1}{2}\,e^{-2t}, & t > 1. \end{cases}}
\]

\textbf{Remark.} Although $y$ is continuous at $t = 1$, the derivative $y'$ is not: on
the left $y'(1^-) = e^{-2}$, while on the right $y'(1^+) = -(e^2-1)e^{-2}$, confirming
the observation made in the problem statement.

\end{respuesta}


\textbf{27. Solve the initial value problem
    \[ y' + p(t)y = 0, \qquad y(0) = 1, \]
    where
    \[ p(t) = \begin{cases} 2, & 0 \le t \le 1, \\ 1, & t > 1. \end{cases} \]}
\begin{respuesta}

The function $p(t)$ has a jump discontinuity at $t_0 = 1$, so we solve separately on
$[0,1]$ and $(1,\infty)$, then match at $t = 1$ by requiring continuity of $y$.

\textbf{Solution on $0 \le t \le 1$.}

Here $p(t) = 2$, so the equation is $y' + 2y = 0$. The integrating factor is
$\mu(t) = e^{2t}$, giving:

\[
\frac{d}{dt}\left(e^{2t}y\right) = 0 \implies e^{2t}y = C_1 \implies y_1 = C_1 e^{-2t}
\]

Applying $y(0) = 1$:

\[
1 = C_1 \implies C_1 = 1
\]

\[
y_1 = e^{-2t}, \qquad 0 \le t \le 1
\]

\textbf{Solution on $t > 1$.}

Here $p(t) = 1$, so the equation is $y' + y = 0$. The integrating factor is
$\mu(t) = e^{t}$, giving:

\[
\frac{d}{dt}\left(e^{t}y\right) = 0 \implies e^{t}y = C_2 \implies y_2 = C_2 e^{-t}
\]

\textbf{Matching at $t = 1$.}

We require $y_1(1) = y_2(1)$:

\[
e^{-2} = C_2 e^{-1} \implies C_2 = e^{-1}
\]

\textbf{Complete solution.}

\[
\boxed{y = \begin{cases} e^{-2t}, & 0 \le t \le 1, \\[6pt] e^{-1}\,e^{-t}, & t > 1. \end{cases}}
\]

which can be written more compactly for $t > 1$ as $y = e^{-(t+1)}$.

\textbf{Remark.} As expected, $y$ is continuous at $t=1$ since both pieces give
$y(1) = e^{-2}$, but $y'$ is not: on the left $y'(1^-) = -2e^{-2}$, while on the
right $y'(1^+) = -e^{-2}$, reflecting the jump in $p(t)$ at that point.

\end{respuesta}









\end{quote}



